\section{结论}
\label{sec:conclusion}
\subsection{四个问题的主要结果}
\subsubsection{安全能力与有限资源运输}
本文面向山区洪涝救援的80个不可拆货箱，以统一地形、载荷、时间、能耗和充电模型贯通四问。在20\%返航余量下，问题一以18架次达到解析下界，A/B/C架次为0/9/9，总能耗59.131296 kWh、累计作业时间546.266929 min；在相同架次数下，混合机型较仅C型方案节能约20.02\%。原组批在安全余量提高至23.08835\%时保持可行与词典序最优，完整参数剖面揭示货箱不可拆分产生的阶梯变化。

问题二联合选择任务、机型、飞机、电池和开始时刻，得到23架次零迟到可行方案，最后返航5693.231489 s，总能耗66.214460 kWh。独立全局有效下界为5433.956428 s，以上界归一化的认证间隙4.5541\%。相同日程的箱号、交付记录和双资源占用一一对应，八架飞机作业占用率95.76\%。充电时间增加30\%时，保持任务与资源顺序的顺延修复可在6108.392219 s完成；超出任务能量预算的情景则由重新组批备选处理。

\subsubsection{连续通信与独立资源配置}
问题三把直连失效转化为完整通信需求区间，联合安排中继位置、高度、服务提供方及资源时序。主方案使用23个运输架次和4个中继架次，联合最后返航5836.969929 s，总能耗68.939416 kWh。中继能源占比约3.90\%，承担约39.57\%的累计运输通信保障；同一悬停兼容物理模型下，关键交接改善使联合工期缩短约87.998 s。附加0.25 dB保障备选的额外工期为29.171 s，12份中继上线延迟修复均通过完整核验。

问题四冻结主联合日程后，15个服务区压缩为3个不可拆任务块。完整枚举得到两组最少缺1架C型运输机、三组缺7件指定型号资源。相对27件共享配置，两组需要28件、三组34件。逐组资源链构造实现了这些峰值，所选两组需求向量对另外两种分区的分量支配，进一步证明固定日程和目标顺序下的选择不随非负库存变化而改变。

\subsection{救援协同执行的决策启示}
\subsubsection{从局部能力转向完整时间轴}
单项任务的最大载荷、最低能耗和最少架次，不能分别替代整个系统的执行评价。山区几何确定任务物理成本，货箱与访问顺序决定交接和载荷过程，飞机与能源恢复决定接续，通信提供方决定关键任务的最早可启动时刻。本文将这些关系放在同一任务和时间体系中，形成可复核的完整救援预案。

\subsubsection{把方案效率、保障等级与资源投入共同呈现}
名义高效率方案、增强保障备选和固定结构恢复分别适合不同执行条件；独立组织方式则进一步改变类型资源需求。本文同时给出方案、边界和调整代价，使决策者能够按任务保障与装备条件选择，而非只追求一个缺乏执行含义的综合分数。统一模型、分层算法、完整结果表与四问压力分析共同构成了从建模、求解到部署的证据链。
